\section{The elliptic construction and its physical end}
\label{so:section}

We prove the existence assertion needed for the filled comparison.  Throughout
this section the dimension is three, so that $k=2$.  We fix the prepared
exterior, $\epsilon$, and then $N>\max(2,2/\epsilon)$.  Constants in this
section may depend on these fixed choices without a specified rate.  The
bounds whose dependence on $N$ matters are precisely those in
Proposition~\ref{sy:coercivity}, Lemma~\ref{es:floor},
and Proposition~\ref{es:distortion}.  In particular, none of the elliptic constants below
is used as a polynomial bound in $N$.
We use the scalar linear elliptic estimates, maximum principles, and
degree theory in \cite[Chapters 6, 8, 9, and 11]{GilbargTrudinger2001}.
The special estimates for the present coupled system are proved below.

\subsection{A scalar estimate with a measurable axial coefficient}

The following estimate is the reason that the trace equation can be treated
before the second unknown has a modulus of continuity.  In its statement,
the matrix $A$ acts on covectors using the metric.

\begin{lemma}[The rough axial estimate]\label{so:rough-gradient}
Let $g$ range over metrics with uniform ellipticity and bounded local
$C^2$ norms on a three-dimensional ball.  Suppose that $z\in C^2$,
$|z|+|dz|\le M$, and
\begin{equation}\label{so:axial-equation}
 A:\nabla^2z=G,\qquad
 A=I-(1-\chi)e\otimes e,\qquad
 e=\frac{\nabla z}{|\nabla z|},\qquad
 0<c\le\chi\le1,\qquad |G|\le M,
\end{equation}
where the axis at a zero of $dz$ may be arbitrary.  No continuity assumption
on $\chi$ is imposed.  There are $\alpha\in(0,1)$ and $C$ depending only
on the displayed bounds and the patch geometry such that on every smaller
patch
\begin{equation}\label{so:rough-estimates}
 [dz]_{C^\alpha}\le C,\qquad
 \int_{B_r(x)}|\nabla^2z|^2\,dV_g\le C r^{1+2\alpha}.
\end{equation}
The same assertion holds up to a smooth face on which $z$ is constant;
balls in the integral are then intersected with the domain.  It is enough
that $z$ be individually $C^1\cap W^{2,s}$ for every finite $s$ and
$C^2$ up to each side of such a reflected face.
\end{lemma}

\begin{proof}
We give the compactness and improvement arguments, including the treatment
of the face.  The identities involving third derivatives are distributional.
For a $C^2$ function and metric their trace-reversed form follows by
$C^2$ smoothing: the fluxes converge locally uniformly, the Hessian
quadratics converge in $L^1$, and the Ricci terms converge locally
uniformly.  The original axial equation is substituted only after this
identity has been established; neither $\chi$ nor $G$ is differentiated
or mollified.  The same argument on each side of a face retains the
normal flux jump.  Thus the proof applies in particular to the
$C^{2,\alpha}$ functions arising in continuation below.

\emph{A Hessian estimate above exponent two.}
Normalize the metric at the center of a small coordinate ball.  The
coordinate principal matrix $a$ in Equation~\ref{so:axial-equation}
satisfies
\[
 |I-a|_{\mathrm F}\le1-c/2
\]
after shrinking the ball by an amount determined by the metric bounds.
For a compactly supported function $Y$ on Euclidean space, the Fourier
transform gives
$\|D^2Y\|_2=\|\Delta Y\|_2$, where the matrix norm is Frobenius.
The Calder\'on--Zygmund estimate at any fixed exponent larger than two
\cite[Theorem 9.9]{GilbargTrudinger2001},
interpolated with this identity, implies that the norm of
$D^2\Delta^{-1}:L^s\longrightarrow L^s$ tends to one as $s\downarrow2$.
Choose $s_0>2$ sufficiently close to two that this norm times $1-c/2$
is less than one.  Writing
$\Delta Y=(I-a):D^2Y+a:D^2Y$ and absorbing gives the global estimate.
Apply it to $\zeta Y$, and use
\[
 \|DY\|_{L^s(B_r)}
 \le \varepsilon r\|D^2Y\|_{L^s(B_r)}
       +C_{\varepsilon}r^{-1}\|Y\|_{L^s(B_r)}.
\]
Using intermediate balls and absorbing successively gives, for
$2\le s\le s_0$,
\begin{equation}\label{so:cordes-local}
 \|D^2Y\|_{L^s(B_{r/2})}
 \le C\bigl(\|a:D^2Y\|_{L^s(B_r)}
                +r^{-2}\|Y\|_{L^s(B_r)}\bigr).
\end{equation}
For clarity, the successive absorption uses radii tending geometrically
from $r/2$ to $r$: choose the interpolation coefficient smaller than the
inverse of the geometric loss to the power $2s$.  The resulting series
converges.  This proves Equation~\ref{so:cordes-local} without a derivative
of $a$.

\emph{The trace-reversed identity.}
In Euclidean coordinates put $H=D^2z$, $q=1-\chi$, and
$W=|Dz|^2/2$.  Since $He= D W/|Dz|$ wherever $Dz\ne0$, direct multiplication
gives the identity, also valid when $Dz=0$,
\begin{equation}\label{so:trace-reversed}
 A DW-GDz=(H-(\operatorname{tr}H)I)Dz,
 \qquad
 \operatorname{div}\bigl((H-(\operatorname{tr}H)I)Dz\bigr)
       =|H|^2-(\operatorname{tr}H)^2.
\end{equation}
Moreover $\operatorname{tr}H=qH(e,e)+G$.  Choose
$(1-c)^2<\theta<1$.  Young's inequality then gives
\[
 (\operatorname{tr}H)^2\le\theta |H|^2+C_cG^2.
\]
If $|Dz|\le1$ and $S=1-|Dz|^2$, we consequently have
\begin{equation}\label{so:deficit-inequality}
 \operatorname{div}(A DS+2GDz)
       \le-\kappa |D^2z|^2+C_cG^2,\qquad \kappa>0.
\end{equation}
This is a divergence inequality with measurable coefficients; it has no
$D\chi$ term.  Intrinsically the last member of
Equation~\ref{so:trace-reversed} gains
$\operatorname{Ric}_g(\nabla z,\nabla z)$.  Passing to coordinates with the
volume density and replacing covariant by ordinary derivatives adds
bounded flux and scalar errors.  On rescaled balls whose metric tends to
the Euclidean metric in $C^1$, whose smooth-side curvature tends to zero,
and whose normalized forcing tends to zero, these errors have the form
\begin{equation}\label{so:small-deficit}
 \operatorname{div}(a_1 DS+B)
 \le-\kappa_1|\nabla^2z|^2+\varepsilon,
 \qquad |B|\le\varepsilon,
\end{equation}
with fixed ellipticity constants and $\varepsilon\to0$.

\emph{A gradient drop or entry into an affine regime.}
Fix $0<r_*<1/32$.  For every $b_0>0$ there exist $k_*\in(r_*,1)$ and
an error tolerance $\varepsilon_*>0$ such that a normalized solution on
$B_1$, with $|\nabla z|\le1$ and errors at most $\varepsilon_*$, satisfies
one of
\begin{equation}\label{so:drop-or-affine}
 \sup_{B_{r_*}}|\nabla z|\le k_*;
 \qquad
 \sup_{B_{r_*}}|z-L|\le b_0r_*,\quad \tfrac12\le|DL|\le2
\end{equation}
for an affine $L$.  To prove this, suppose the assertion fails, let the
errors tend to zero, and let the gradient supremum on $B_{r_*}$ tend to
one.  The scalar weak Harnack inequality
\cite[Theorem 8.18]{GilbargTrudinger2001}, applied to the nonnegative
supersolution $S$ in Equation~\ref{so:small-deficit}, gives
\[
 \left(\frac{1}{|B_{1/3}|}\int_{B_{1/3}} S^p\right)^{1/p}
 \le C\bigl(\inf_{B_{r_*}}S+\varepsilon\bigr)\longrightarrow0
\]
for a fixed $p>0$.  Equation~\ref{so:cordes-local} gives a uniform local
$L^2$ bound for $DS$.
Here is the extra truncation step needed to retain the negative Hessian
term.  Testing the supersolution inequality by
$\zeta^2(b-S)_+$ yields
\begin{equation}\label{so:lower-truncation}
 \int_{\{S<b\}}\zeta^2|DS|^2
 \le Cb^2\int|D\zeta|^2+C\|B\|_\infty^2
       +C\varepsilon b.
\end{equation}
Indeed the derivative of $(b-S)_+$ contributes the negative of the
elliptic energy on $\{S<b\}$; the cutoff and $B$ terms are absorbed by
Young's inequality.  Since $S\to0$ in measure, split $DS$ into the sets
$\{S<b\}$ and $\{S\ge b\}$.  The first has limiting $L^1$ norm at most
$Cb$ by Equation~\ref{so:lower-truncation}; the second has vanishing
$L^1$ norm by the uniform $L^2$ bound and its vanishing measure.  Thus
$DS\to0$ locally in $L^1$.  Testing Equation~\ref{so:small-deficit} with
a fixed nonnegative cutoff now gives $D^2z\to0$ locally in $L^2$.
After subtraction of a constant, the gradient bound gives uniform
compactness, and the limit is affine.  Its slope has length one because
$S\to0$ in measure.  This contradicts failure of the second alternative
in Equation~\ref{so:drop-or-affine}.

\emph{Improvement near a nonzero affine function.}
Suppose on $B_1$ that $|z-L|\le b\le b_0$, $1/4\le|DL|\le4$, and the
gradient is bounded by a fixed constant.  Suppose also that metric
oscillation, connection errors, and forcing are at most $b\delta$.
For $Y=(z-L)/b$, Equation~\ref{so:cordes-local} gives a uniform
$W^{2,s_0}$ bound on smaller balls; its source is bounded by $C\delta$.
Put $e_0=DL/|DL|$.  The directions satisfy
\[
 |e\otimes e-e_0\otimes e_0|
       \le C\min\{1,b|DY|\}+o(1).
\]
Hence they converge in measure to the fixed axis as $b_0,\delta\to0$.
H\"older's inequality, with exponent
$2s_0/(s_0-2)$ for this bounded coefficient error, proves
\begin{equation}\label{so:frozen-residual}
 \|R\|_{L^2(B_{3/4})}\longrightarrow0,
 \qquad R=(\Delta_{e_0^\perp}+\chi\partial_{e_0}^2)Y.
\end{equation}
Only the axis is frozen; $\chi$ remains an arbitrary measurable function
of position.

There is a zero-Dirichlet correction $U\in W^{2,2}(B_{3/4})\cap
W^{1,2}_0(B_{3/4})$ with
\[
 (\Delta_{e_0^\perp}+\chi\partial_{e_0}^2)U=R,
 \qquad \|U\|_{W^{2,2}}\le C_c\|R\|_2.
\]
Indeed solve $U=\Delta_D^{-1}((1-\chi)U_{e_0e_0}+R)$ by contraction
in the Hessian norm.  On a convex ball integration by parts gives
\[
 \int |D^2U|^2\le\int(\Delta U)^2;
\]
the difference is the boundary mean curvature times $(\partial_\nu U)^2$,
which is nonnegative.  Thus the contraction factor is at most $1-c$.
Poincar\'e's inequality controls the lower derivatives.  In dimension
three the Sobolev embedding $W^{2,2}\hookrightarrow C^{0,1/2}$ gives
$\|U\|_\infty\to0$.  This is the dimension-specific uniform correction
required here.

Write $Y_0=Y-U$, and rotate so that $e_0$ is the third coordinate axis.
Then $v=\partial_3Y_0\in W^{1,2}$ satisfies
\begin{equation}\label{so:axial-derivative}
 \partial_1^2v+\partial_2^2v+\partial_3(\chi\partial_3v)=0
\end{equation}
distributionally: it is exactly the distributional derivative of
$\Delta_{e_0^\perp}Y_0+\chi\partial_3^2Y_0=0$.
The scalar divergence-form H\"older estimate
\cite[Chapter 8]{GilbargTrudinger2001} and its energy inequality
therefore give, on smaller balls,
\[
 [v]_{C^{\alpha_1}}\le C,\qquad
 \int_{B_r}|Dv|^2\le Cr^{1+2\alpha_1}
\]
for some $\alpha_1\in(0,1)$.  Since
$\Delta Y_0=(1-\chi)\partial_3v=:F_0$, Cauchy--Schwarz gives
\begin{equation}\label{so:newton-mass}
 \int_{B_r}|F_0|\le Cr^{2+\alpha_1}.
\end{equation}
This also controls all the first derivatives of $Y_0$.  To see the
point explicitly, localize $F_0$ to a slightly larger ball and take its
Newton potential $P$.  In the difference $DP(x)-DP(y)$, with
$h=|x-y|$, the two near regions have size bounded by
$C\sum_{j\ge0}(2^{-j}h)^{\alpha_1}\le Ch^{\alpha_1}$, by
Equation~\ref{so:newton-mass}.  On the annuli of radius $r\ge2h$,
the derivative of the gradient kernel is bounded by $Cr^{-3}$, giving
$Ch\sum_{r=2^jh}r^{\alpha_1-1}\le Ch^{\alpha_1}$.
The remainder $Y_0-P$ is harmonic.  Its local norms are bounded by the
already established $W^{2,2}$ bound and the potential bound.  Consequently
$DY_0$ is uniformly H\"older on an interior ball.

Choose $0<\alpha_2<\alpha_1$, then a sufficiently small fixed
$\lambda\in(0,1/8)$ so that the Taylor error of $Y_0$ is at most
$\frac12\lambda^{1+\alpha_2}$.  Finally choose $b_0$ and $\delta$ so
that the uniform correction is smaller than the other half.  We obtain
an affine $L'$ with
\begin{equation}\label{so:affine-improvement}
 \sup_{B_\lambda}|z-L'|\le b\lambda^{1+\alpha_2},
 \qquad |DL'-DL|\le Cb.
\end{equation}

\emph{Iteration and the boundary interface.}
In the improvement statement $b$ is an allowed upper bound for the
error, and need not equal the actual error.  Fix $b_0,\delta,\lambda$
as above, decreasing $b_0$ so that
$Cb_0/(1-\lambda^{\alpha_2})<1/4$, and then choose $k_*,\varepsilon_*$
in Equation~\ref{so:drop-or-affine}.  Initially shrink to a radius $R_0$
and normalize the gradient by a fixed bound $\Lambda_0$ so that all errors
are at most $\min(\varepsilon_*,b_0\delta)$.  After $j$ drops the
physical radius and gradient normalization are
$R_j=R_0r_*^j$ and $\Lambda_j=\Lambda_0k_*^j$.  In particular the normalized
forcing is $R_jG/\Lambda_j$, and is multiplied by $r_*/k_*<1$ at each drop.
Metric and connection errors decrease at least by $r_*$ and curvature
errors by $r_*^2$.  At the first affine entry the physical radius is
$R_e=r_*R_j$, with the same $\Lambda_j$.  Rescale that ball and apply
Equation~\ref{so:affine-improvement} with the prescribed upper bound
$b=b_0$, even if the actual affine error is smaller.  On iteration $i$
use $b_i=b_0\lambda^{i\alpha_2}$.  The forcing divided by $b_i$ is
at most $\delta\lambda^{i(1-\alpha_2)}$; the same comparison controls
the connection errors.  The slope increments are summable by the
choice of $b_0$, so the axis remains separated from zero.  Only
constants are subtracted in the rescaled original equation, so its
axis is always the actual gradient axis.  These two regimes give
uniform affine approximation of order $r^{1+\alpha}$, including the
transition scale, for any sufficiently small
\[
 0<\alpha\le
 \min\{\alpha_2,\log k_*/\log r_*\}.
\]
The affine approximation criterion for H\"older gradients can also be
seen directly: affine slopes at consecutive dyadic scales differ by
$Cr^\alpha$, hence converge; comparing overlapping balls gives the
same modulus between two centers.  Thus the first estimate in
Equation~\ref{so:rough-estimates} follows.  Subtract the first jet in
Equation~\ref{so:cordes-local}.  Its error is $O(r^{1+\alpha})$;
connection terms are bounded.  Squaring the resulting $L^2$ estimate
proves the second estimate.

For a face use Gaussian coordinates $(x',x_3)$, subtract its constant
value, reflect $z$ oddly and $g$ evenly.  The reflected function is
$C^1\cap W^{2,s}$, the metric is Lipschitz, and each smooth side has
bounded curvature.  The coordinate Hessian estimate applies unchanged.
In the trace-reversed identity the only additional term is the jump of
the normal flux at $x_3=0$.  Since the tangential gradient vanishes there,
\[
 \bigl((\nabla^2z-(\Delta z)g)\nabla z\bigr)\cdot\nu
       =-(\partial_\nu z)\operatorname{tr}_{T}\nabla^2z.
\]
The constant boundary value gives
$|\operatorname{tr}_{T}\nabla^2z|\le C|\mathrm{II}|\,|dz|$.
Thus the jump of the density-weighted flux is a bounded plane density
$J(x')$, of size $C|\mathrm{II}|\,|dz|^2$.  It can be cancelled exactly
by adding the coordinate vector
$-J(x')\mathbf1_{\{x_3>0\}}\partial_3$ to the flux; its divergence
is $-J\delta_{\{x_3=0\}}$, and no tangential derivative of $J$ is used.
Under a spatial scale $R_j$ and gradient normalization $\Lambda_j$, the
normal flux jump becomes $R_jJ/\Lambda_j^2$.  Since the current gradient
bound is $|dz|\le\Lambda_j$, its size is at most $CR_j|\mathrm{II}|$,
also after any number of gradient drops.  The smooth-side Ricci
error is at most $CR_j^2|\operatorname{Ric}|$.  Both therefore tend
to zero independently of the shrinking gradient normalization.
Hence Equation~\ref{so:small-deficit}, the compactness
alternative, and the affine improvement all remain valid across the
face.  This proves the boundary assertion.
\end{proof}

\subsection{Regularity of the two unknowns}

\begin{lemma}[A small potential estimate]\label{so:potential}
Let $M$ be a symmetric bounded measurable uniformly elliptic matrix on
$B_r\subset\mathbb R^3$.  If $|B|\le C_0$ and a nonnegative finite
measure $\mu$ satisfies
$\mu(B_s(x))\le C_0s^{1+\gamma}$ for some $0<\gamma<1$, then the
zero-Dirichlet solution of
\[
 -\operatorname{div}(MDv)=\operatorname{div}B+\mu
\]
satisfies $\|v\|_\infty\le C(r+r^\gamma)$.  The assertion applies
uniformly to smooth approximations of the data and hence to their
weak limits.
\end{lemma}

\begin{proof}
For the bounded flux part, scale to the unit ball.  Testing by
$(v-k)_+$ gives
$\int_{\{v>k\}}|Dv|^2\le C\|B\|_\infty^2|\{v>k\}|$.
Sobolev's inequality and the level-set iteration give
$\|v\|_\infty\le C\|B\|_\infty$ on the unit ball; the same
argument for $-v$ gives the other sign.  Rescaling gives $CrC_0$.
For the measure part the scalar Dirichlet Green function has the
bound $0\le G(x,y)\le C|x-y|^{-1}$ for bounded measurable symmetric
uniformly elliptic coefficients.  This follows from the whole-space
bound in \cite[Theorem 7.1 and Equation (7.9)]{LittmanStampacchiaWeinberger1963}
by domain comparison.  The coefficients may be extended elliptically
outside the ball; symmetry is preserved.  Integrating on dyadic balls centered at $x$
gives
\[
 \int G(x,y)\,d\mu(y)
 \le C\sum_{j\ge0}(2^{-j}r)^{-1}
                      (2^{-j}r)^{1+\gamma}
 \le Cr^\gamma.
\]
For completeness the measure belongs to $H^{-1}$: its Newton energy
is bounded by the displayed potential bound times $\mu(B_r)$, and
the Newton energy identity bounds its action on $H^1_0$.  Mollifications
preserve the mass-growth bound uniformly.  Indeed, at scales
$s<\varepsilon$ their density is bounded by $C\varepsilon^{\gamma-2}$,
so their mass is at most
$Cs^3\varepsilon^{\gamma-2}\le Cs^{1+\gamma}$; at larger scales use
the original bound on a ball enlarged by $\varepsilon$.
The bounded divergence-data estimate can first be applied to smooth
matrices and then passed to the weak solution by its energy bound.
The measure approximations have uniformly bounded $H^{-1}$ norm and
uniformly bounded potentials, so weak $H^1_0$ compactness and weak-star
$L^\infty$ compactness give the same estimate for their limit.
This proves the lemma.
\end{proof}

\begin{proposition}[Coupled estimates]\label{so:coupled-regularity}
On a fixed truncation, every smooth solution of the homotopy in
Equation~\ref{sy:homotopy} satisfying $t\ge-\epsilon$ has local smooth
bounds depending only on the fixed data and $N$.  The bounds hold on
uniformly sized interior patches and on outer Dirichlet half-patches,
independently of a sufficiently large truncation radius.  In particular
there is a fixed exponent $\alpha_0>0$ and a uniform local
$C^{1,\alpha_0}$ bound for $f$ and $C^{0,\alpha_0}$ bound for $Z$.
\end{proposition}

\begin{proof}
The bounds in Proposition~\ref{es:distortion} give uniform ellipticity
of $A_\chi$ and boundedness of every smooth function of $x,Z,df$ that
occurs in the equations.  The trace equation, divided by
$lD^{-1/2}$, is
\begin{equation}\label{so:normalized-trace}
 A_\chi:\nabla^2 f
       =\frac{\sqrt D}{l}\bigl(\eta f-
                                  \operatorname{tr}_{A_\chi}K_a\bigr).
\end{equation}
Its right side is bounded.  Lemma~\ref{so:rough-gradient} therefore
bounds $df$ in $C^\alpha$ and gives
$\int_{B_r}|\nabla^2f|^2\le Cr^{1+2\alpha}$.
The exact differential identity in Equation~\ref{sy:differential}
gives
\begin{equation}\label{so:t-derivative-bound}
 |dt|\le C(|dZ|+|\nabla^2 f|+1).
\end{equation}
Write the second equation as
\[
 \operatorname{div}(M dZ+J)=F_1,\qquad
 M=2uA_\chi,\quad J=uA_\chi K_a(w,\cdot).
\]
Here $M$ is symmetric uniformly elliptic and $J$ is bounded.  The
quadratic formula for $\mathcal T$ and
Equation~\ref{so:t-derivative-bound} imply
\begin{equation}\label{so:natural-growth}
 |F_1|\le C(1+|dZ|^2)+C|\nabla^2f|^2.
\end{equation}
All volume densities and connection terms can be incorporated into
$M,J,F_1$ without changing these assertions.  Set
$d\mu=(1+|\nabla^2f|^2)\,dx$.  For a fixed
$0<\gamma<\min(1,2\alpha)$ this measure satisfies
$\mu(B_r)\le Cr^{1+\gamma}$.

Choose a sufficiently large fixed $L$ and put $Q_\pm=e^{\pm LZ}$.
For either sign the chain rule gives
\[
 -\operatorname{div}(M dQ_\pm)
 =\operatorname{div}(Q_\pm'J)-Q_\pm'F_1
       -Q_\pm''\bigl(M(dZ,dZ)+J\cdot dZ\bigr).
\]
The range of $Z$ is bounded.  Ellipticity and Young's inequality show
that the last term absorbs the quadratic term in
Equation~\ref{so:natural-growth} if $L$ is large.  Thus
\begin{equation}\label{so:exponential-subsolutions}
 -\operatorname{div}(M dQ_\pm)
       \le\operatorname{div}B_\pm+C\mu,\qquad |B_\pm|\le C.
\end{equation}
By Lemma~\ref{so:potential}, subtracting the zero-Dirichlet potential
of the right side on $B_r$ changes $Q_\pm$ by at most
$C(r+r^\gamma)$ and makes it a subsolution.  The corrected supremum
deficit is nonnegative and is a supersolution.  Apply weak Harnack to
that deficit.  On at least half of $B_{r/2}$, either $Z$ is below the
midpoint of its range on $B_r$, or it is above that midpoint.  In the
first case the $Q_+$ deficit, and in the second the $Q_-$ deficit, is
at least a fixed multiple of $\operatorname{osc}_{B_r}Z$, up to the
potential error.  Weak Harnack then decreases one end of the range on
$B_{r/4}$ by a fixed fraction.  More precisely,
\[
 \operatorname{osc}_{B_{r/4}}Z
 \le\theta\operatorname{osc}_{B_r}Z+C(r+r^\gamma),
 \qquad 0<\theta<1.
\]
Iteration proves a uniform H\"older modulus for $Z$.
At an outer face both unknowns have zero trace.  In flattened
coordinates let $P=\operatorname{diag}(1,1,-1)$.  After incorporating
the volume density, reflect by
$Z^-=-Z^+$, $M^-=PM^+P$, $J^-=-PJ^+$, and $F_1^-=-F_1^+$,
where opposite-side values are evaluated at reflected points.
The reflected total flux is $-P(M^+dZ^++J^+)$, so its normal component
has matching traces and creates no plane source.  For the principal
part this also follows from the vanishing tangential gradient on the
constant-data face.  In the actual application $K_a=0$ near the outer
face, hence $J=0$ there.  The equation and its growth estimate hold
across the face, and the same oscillation argument applies there.

The coefficients and right side of Equation~\ref{so:normalized-trace}
are now H\"older.  Interior and Dirichlet Schauder estimates give
$C^{2,\alpha'}$ bounds for $f$, with some $\alpha'>0$.  Expanding the
second equation, all derivatives of its coefficients are derivatives
of smooth functions of $x,Z,df$.  Consequently it has the form
\begin{equation}\label{so:Z-nondivergence}
 M^{ij}\partial_{ij}Z=H(x,Z,dZ),\qquad
 |H(x,Z,dZ)|\le C(1+|dZ|^2),
\end{equation}
with a uniformly H\"older principal matrix.  No bound for $dZ$ has
been used to reach this point.

Here is the absorption which supplies that bound.  For $s>3$ the local
$W^{2,s}$ estimate for a H\"older principal matrix has a uniformly
bounded leading constant on sufficiently small patches.  On balls
or constant-Dirichlet half-balls the scaled derivative interpolation
inequality is
\begin{equation}\label{so:gradient-interpolation}
 \|dZ\|_{L^{2s}(Q)}^2
 \le C\operatorname{osc}_{Q'}Z\,
                      \|D^2Z\|_{L^s(Q')}
        +C r^{3/s-2}(\operatorname{osc}_{Q'}Z)^2,
\end{equation}
where $Q\Subset Q'$ are concentric patches of comparable radius $r$.
It follows by applying the $L^\infty$--$W^{2,s}$ interpolation
inequality to a cutoff times $Z-c$; take $c$ between its maximum and
minimum.  For half-balls subtract the zero boundary value and reflect,
or use a boundary cutoff in flattened coordinates.
The H\"older modulus makes the oscillation coefficient as small as
needed by fixing a sufficiently small patch radius.  To justify
absorption locally without any bound depending on the number of
patches, let $X$ be the supremum of their $L^s$ Hessian norms.  Each
enlarged patch is covered by a fixed number of the smaller patches.
The local estimate and Equation~\ref{so:gradient-interpolation} give
$X\le\frac12X+C$.  This supremum is individually finite on every
smooth compact truncation.  Thus $X\le2C$, independently of its
radius.  Sobolev embedding gives a H\"older gradient for $Z$.
Schauder estimates applied to both equations now give second derivative
bounds; differentiating gives bounds of every fixed order.  The
background geometry, end symbols, and outer sphere charts have
uniform bounds at the chosen patch scale.  The constants are
therefore independent of the sufficiently large truncation radius.
\end{proof}

\subsection{The trace solve for arbitrary trials}

\begin{proposition}[The unrestricted trial trace problem]\label{so:trial-trace}
Fix a smooth compact truncation $\Omega_{N,R}$, $a\in[0,1]$,
and $\alpha\in(0,1)$.
For every $Z\in C^{1,\alpha}(\overline{\Omega_{N,R}})$, including trials
for which the implicit $t$ is below $-\epsilon$, there is a unique
solution $f\in C^{3,\alpha}(\overline{\Omega_{N,R}})$ of
\[
 \operatorname{tr}_{A_\chi}(K_a+H^f)=\eta f,
 \qquad f|_{\partial\Omega_{N,R}}=0.
\]
The solution depends continuously on $(a,Z)$ in $C^{2,\alpha}$; bounded
sets of trials have uniform bounds sufficient for this assertion.
For $a=0$ the solution is identically zero.
\end{proposition}

\begin{proof}
All constants here may depend on a bound for the trial's
$C^{1,\alpha}$ norm and on this fixed domain.  The global implicit
relation in Equation~\ref{sy:implicit-t} makes $t$, $l$, $D$, and
$A_\chi$ smooth functions of $(x,Z,df)$, even at $df=0$.
The normalized equation is
\[
 A_\chi:\nabla^2f-qf+\frac{\sqrt D}{l}
                           \operatorname{tr}_{A_\chi}K_a=0,
 \qquad q=\eta\frac{\sqrt D}{l}>0.
\]
Continue it from $\Delta_gf=f$ by using, for $\sigma\in[0,1]$,
\begin{equation}\label{so:trace-continuation}
 A_\sigma:\nabla^2f-q_\sigma f
       +\sigma\frac{\sqrt D}{l}
                             \operatorname{tr}_{A_\chi}K_a=0,
 \quad
 A_\sigma=(1-\sigma)I+\sigma A_\chi,
 \quad q_\sigma=1-\sigma+\sigma q.
\end{equation}
The coefficient of $f$ is strictly negative in this convention.
At a positive maximum $df=0$, so $t=Z$, $D=1$, and $A_\chi=I$.
The maximum principle gives
\[
 |f|\le H:=\eta^{-1}\|\operatorname{tr}_gK\|_\infty
\]
for the entire continuation.  If $K=0$ it gives $f=0$.

We prove a gradient bound before using uniform ellipticity.  Write
$s=|df|$.  When $s\to\infty$ with $Z$ in a bounded interval, the
implicit equation forces $t\to-\infty$.  There $p=N$ and $l=e^{Nt}$,
and the defining equation gives
\begin{equation}\label{so:large-trial-slope}
 t=-\frac{\log s}{N+2}+O(1),\qquad
 d\asymp s^{-4/(N+2)},\qquad
 \chi\ge\frac{c}{1+s},\qquad
 \frac{\sqrt D}{l}\le C(1+s).
\end{equation}
These estimates are uniform for bounded trial values; on a bounded
$s$ interval they follow from smoothness and positivity.  The lower
bound for $\chi$ uses $N>2$.  Notice that it does not use the floor.
The axial eigenvalue $\chi_\sigma=1-\sigma+\sigma\chi$ obeys the
same lower bound, while both transverse eigenvalues equal one.
Since $|f|\le H$, Equation~\ref{so:trace-continuation} has normalized
source of absolute value at most $C(1+s)$.

Take a short fixed collar of the boundary and a smooth extension of
the background metric across it.  Choose a concave increasing modulus
\[
 \psi(d)=K_0^{-1}\log(1+K_0Md),\qquad
 \psi''=-K_0(\psi')^2.
\]
First choose $K_0$ sufficiently large depending on the source and
geometry bounds.  Then choose a sufficiently short $d_0$ and a
sufficiently large $M$ so that $\psi'\ge1$ on $[0,d_0]$ and
$\psi(d_0)>2H$.  Such choices are compatible: one can first make
$K_0d_0<1/2$ and then let $M\to\infty$.
For the distance $d$ to the boundary, the Hessian in its gradient
axis is $\psi''$; the transverse trace is bounded by $C\psi'$.
Equation~\ref{so:large-trial-slope} gives
\[
 \chi_\sigma\psi''+C\psi'
 \le-cK_0\frac{(\psi')^2}{1+\psi'}+C\psi'
 \le(-cK_0/2+C)\psi'.
\]
For large $K_0$ this is strictly more negative than the possible source.
The negative barrier has the opposite strict inequality.  At a first
contact the gradients agree, and hence all the coefficients depending
on the gradient agree.  The strict sign and positivity of $q_\sigma$
therefore exclude contact.  The boundary values and the
choice $\psi(d_0)>H$ give
\begin{equation}\label{so:trace-face-modulus}
 |f(x)|\le\psi(\operatorname{dist}(x,\partial\Omega_{N,R}))
 \quad\hbox{in the collar}.
\end{equation}

For the interior consider
$f(x)-f(y)-\psi(\operatorname{dist}(x,y))$ for distances at most $d_0$,
using the short geodesic distance in the smooth extended metric.
The value is nonpositive on the diagonal, at distance $d_0$, and
whenever one endpoint lies on the boundary, the last assertion by
Equation~\ref{so:trace-face-modulus} and monotonicity of $\psi$.
At a hypothetical positive interior maximum the endpoint gradients
are parallel along the short geodesic and have the same magnitude
$\psi'$.  Write $e_x,e_y$ for that common transported direction.
Varying just the first endpoint in its axial direction and then just
the second endpoint gives, respectively,
\[
 \nabla^2 f_x(e_x,e_x)\le\psi'',\qquad
 \nabla^2 f_y(e_y,e_y)\ge-\psi''.
\]
Paired parallel transverse endpoint variations, summed over the two
transverse directions, give
\[
 \operatorname{tr}_{e_x^\perp}\nabla^2f_x
 -\operatorname{tr}_{e_y^\perp}\nabla^2f_y
 \le C\operatorname{dist}(x,y)\psi'.
\]
This last bound is the second variation of the short geodesic length;
the endpoint Jacobi fields have equal parallel initial and final
values, and their index forms are bounded by curvature times its
length.  Thus the difference of the principal operators is at most
\begin{equation}\label{so:different-axial-values}
 (\chi_\sigma(x)+\chi_\sigma(y))\psi''+Cd_0\psi'
 \le(-cK_0+Cd_0)\psi'.
\end{equation}
The right sides of the two equations have difference bounded below
by $-2C(1+\psi')$.  Increasing $K_0$ gives a contradiction.  Crucially,
Equation~\ref{so:different-axial-values} uses the two axial coefficients
separately; no modulus of $\chi(x)-\chi(y)$ is assumed.  We conclude
$|f(x)-f(y)|\le\psi(\operatorname{dist}(x,y))$ for short distances,
and hence $|df|\le M$.

Uniform ellipticity now follows from this gradient bound.  The matrix
$A_\sigma$ still has precisely the structure of
Lemma~\ref{so:rough-gradient}, and its right side is bounded.  That
lemma gives a H\"older gradient.  Schauder estimates then give
$C^{2,\beta}$ bounds for some $\beta>0$.  At this stage $df$ is
Lipschitz, and the prescribed $Z$ is $C^{1,\alpha}$, so the coefficients
and source of Equation~\ref{so:trace-continuation} are $C^\alpha$.
Schauder estimates give $C^{2,\alpha}$ bounds.  Differentiating the
equation and applying the same estimates, with the boundary estimates
for tangential derivatives followed by the equation for the normal
second derivative, gives $C^{3,\alpha}$ bounds.

The linearization in $f$ is uniformly elliptic with H\"older
coefficients, bounded first-order terms, and zero-order term
$-q_\sigma<0$.  Derivatives of the gradient-dependent coefficients
produce only first-order terms, because those coefficients do not
depend on the value of $f$.  The Dirichlet maximum principle and linear
Schauder solvability make this linearization an isomorphism
$C^{2,\alpha}_0\to C^{0,\alpha}$.  The implicit function theorem gives
openness of the set of continuation parameters.  The uniform
$C^{3,\alpha}$ bounds give compactness in $C^{2,\alpha}$ and hence
closedness.  At $\sigma=0$ the unique solution is zero, so the set is
all of $[0,1]$.
If two solutions of the final equation differed, a positive maximum
of their difference would have matching gradients and an ordered
Hessian, while their common $q$ is positive.
This contradicts the equations.  Thus the solution is unique.
Finally the same implicit function theorem, now with parameters
$(a,Z)$, gives the asserted continuous dependence.  Its local branches
agree globally by uniqueness.
\end{proof}

\subsection{The compact map and the floor}

\begin{theorem}[Filled solver and normalized end]\label{so:filled-solver}
For the choices in Section~\ref{sy:section},
Lemma~\ref{es:floor}, and Proposition~\ref{es:distortion}, every sufficiently
large $N$ admits a smooth solution $(f,Z)$ of Equation~\ref{sy:system}
on the complete filled manifold $\Omega_N$.  It satisfies
\[
 -\epsilon<t\le Z\le\mathcal P_N,\qquad
 \ell\le l\le e,\qquad |f|+|df|\le\exp(\mathcal P_N).
\]
The height estimates of Lemma~\ref{es:height} also hold.  On the physical
end, for each fixed integer $j\ge0$ there are $c_j,C_j>0$, allowed to
depend on $N$, such that
\begin{equation}\label{so:physical-end}
 f=O_j(e^{-c_jr}),\qquad
 t,Z=O_j(r^{-1}),\qquad
 \Delta_g t=O(r^{-3-\beta}).
\end{equation}
In particular both end constants are zero.  The solution is obtained
by exhausting the outer Dirichlet radius with $N$ fixed.  For the
resulting solution $u\Xi$ is absolutely integrable, $u\to1$, and the
penalty vanishes sufficiently far out on the end.
\end{theorem}

\begin{proof}
Fix $R\ge R_{\min}(N)$ as in Lemma~\ref{es:height}.  Parametrize the
homotopy by $\lambda\in[0,2]$:
\[
 (a(\lambda),b(\lambda))=
 \begin{cases}(1-\lambda,1),&0\le\lambda\le1,\\
               (0,2-\lambda),&1\le\lambda\le2.
 \end{cases}
\]
All quantities in the trace and scalar identities use $K_a=aK$;
the divergence equation is $\operatorname{div}V=b\,u\Xi$.
Let $X=C^{1,\alpha}_0(\overline{\Omega_{N,R}})$ for a fixed
$\alpha\in(0,1)$.  For a trial $Z\in X$ solve the trace equation by
Proposition~\ref{so:trial-trace}, and evaluate $t,w,u,A_\chi,\mathcal T$
and $\Xi$ at this pair.  Define $T_\lambda Z=Y$ by the linear
zero-Dirichlet equation
\begin{equation}\label{so:compact-map}
 \operatorname{div}(2uA_\chi dY)
       =b\,u\Xi-\operatorname{div}(uA_\chi K_a(w,\cdot)).
\end{equation}
Only the gradient of the new unknown $Y$ is left unfrozen.
On a bounded set of trials, Proposition~\ref{so:trial-trace} bounds
$f$ in $C^{2,\alpha}$, and all implicit coefficients have bounded
$C^{1,\alpha}$ norms.  The scalar right side has bounded
$C^{0,\alpha}$ norm, since $dt$ is computed from
Equation~\ref{sy:differential}.  Ellipticity is uniform on that bounded
set by the trial gradient estimate, even when $t<-\epsilon$.
Linear Dirichlet Schauder theory gives a uniform $C^{2,\alpha}$ bound
for $Y$.  Thus $(\lambda,Z)\mapsto T_\lambda Z$ is continuous and
compact into $X$.  At a fixed point the linear solve first gives
$Z\in C^{2,\alpha}$, whereas the trace solution is already
$f\in C^{3,\alpha}$.  Thus $\Xi\in C^{1,\alpha}$ and the principal
coefficient and drift flux are $C^{2,\alpha}$.  Schauder estimates give
$Z\in C^{3,\alpha}$, and the trace equation gives
$f\in C^{4,\alpha}$.  Repeating gives smoothness.  In particular the
coupled a priori estimates apply to individually smooth fixed points;
their eventual bounds for all fixed derivative orders control the
chosen $X$ norm, regardless of its initially prescribed exponent
$\alpha$.

Choose a radius $M_X$ larger than the $X$ norm of every above-floor
fixed point, which is possible by Proposition~\ref{so:coupled-regularity}
and the fixed truncation.  Consider the relatively open bounded set
\begin{equation}\label{so:moving-domain}
 \mathcal U=\{(\lambda,Z):\|Z\|_X<M_X,
       \ \min_{\overline{\Omega_{N,R}}}t[Z,df[a(\lambda),Z]]
                                                >-\epsilon\}.
\end{equation}
Its openness follows from the continuous dependence in
Proposition~\ref{so:trial-trace}.  A fixed point cannot occur at its
norm boundary by the choice of $M_X$.  At its other boundary it would
have minimum $t=-\epsilon$ and would be above the closed floor.
Lemma~\ref{es:height} excludes such a minimum on the outer boundary;
Lemma~\ref{es:floor} excludes it in the interior.  There is no
inner boundary in the filled manifold.

For completeness the variation of the domain in
Equation~\ref{so:moving-domain} presents no degree obstruction.
The images of the bounded parameter-trial cylinder under $T$ have
compact closure.  Therefore any convergent sequence of parameters
and associated boundary fixed points would converge to a boundary
fixed point, already excluded.  Around any fixed parameter the compact
set of its fixed points has an open neighborhood whose closure lies
inside its slice of $\mathcal U$.  By compactness the same neighborhood
contains every fixed point in the relevant slices for all sufficiently
nearby parameters; otherwise a contrary sequence converges to an
excluded fixed point.  Excision and the usual fixed-domain homotopy
invariance of Leray--Schauder degree show that the degrees on the
slices are locally constant.  A finite interval covering makes them
constant along the full path.
At $\lambda=2$ we have $K_a=0$, $f=0$, and $b=0$.  The linear solve
in Equation~\ref{so:compact-map} is therefore identically zero for
every trial.  Its sole fixed point is $Z=0$, with $t=0>-\epsilon$,
and the degree of $I-T_2$ on the terminal slice is one.  Hence the
initial slice has degree one and contains a fixed point.  This proves
existence on every sufficiently large truncation, with the estimates
already established.

We next verify the end before taking the exhaustion limit.  All
bounds in the remainder of the proof are uniform in the truncation
radius at this fixed $N$.  Outside a fixed sphere $K=0$, so
Equation~\ref{so:normalized-trace} is the homogeneous proper equation
\begin{equation}\label{so:end-trace}
 A_\chi:\nabla^2 f-q(x)f=0,\qquad
 q(x)=\eta\frac{\sqrt D}{l}\ge\eta/e>0.
\end{equation}
The eigenvalues of $A_\chi$ stay in a fixed positive interval.
For a sufficiently small fixed $c>0$, the functions
$C e^{-c(r-r_0)}$ have
$(A_\chi:\nabla^2-q)C e^{-c(r-r_0)}<0$ for $r_0$ sufficiently large.
Indeed their Hessian is bounded by
$C(c^2+c/r)e^{-c(r-r_0)}$, whereas the zero-order term is at least
$(\eta/e)C e^{-c(r-r_0)}$.  Choose $c$ first, then $r_0$.
The negatives give the opposite inequality.  Choose $C$ to dominate
the uniformly bounded values on $r=r_0$; the outer Dirichlet data
are zero.  Comparison gives $|f|\le Ce^{-cr}$ on every truncated end.
The local estimates in Proposition~\ref{so:coupled-regularity} give
uniform smooth coefficients for Equation~\ref{so:end-trace}; its
interior and boundary estimates then give exponential decay of every
fixed derivative of $f$.  Alternatively, interpolation between this
exponential height bound and the uniform higher derivative bounds
gives the same conclusion, with a possibly smaller positive exponent
for each order.  The implicit relation implies that $t-Z$ and its
fixed derivatives are exponentially small.

On the bounded fixed-$N$ range of $Z$, define the smooth increasing
function $\Phi$ by
\begin{equation}\label{so:poisson-change}
 \Phi(0)=0,\qquad \Phi'(0)=1,\qquad
 \frac{\Phi''(z)}{\Phi'(z)}
      =1+p(z)-\delta_0s_{p(z)},\qquad s_{p}=p+\tfrac12.
\end{equation}
Its derivative is bounded above and away from zero on that range.
Putting $f=K=0$ in Lemma~\ref{sy:zero-trace-reduction} gives exactly
\[
 2\Delta_g Z+2\bigl(1+p(Z)-\delta_0s_{p(Z)}\bigr)|dZ|^2
       =\rho-\vartheta(N(Z+\epsilon))C_N\rho_0.
\]
The exponential bounds just proved and the uniform local derivative
bounds show that the actual end equation differs from this one by
$O_j(e^{-c_jr})$ for every fixed order.  Thus
\begin{equation}\label{so:transformed-poisson}
 |\Delta_g\Phi(Z)|\le C r^{-3-\beta},\qquad
 \Phi(Z)|_{r=R}=0.
\end{equation}
This scalar Poisson bound is independent of the radius $R$.

Here are explicit barriers fixing the end constant.  Put
$\gamma=\beta/2$ and, after increasing $r_0$, let
$W(r)=r^{-1}-r^{-1-\gamma}>0$.  The prepared metric has
$g-\delta=O_j(r^{-1})$.  Since
$\Delta_\delta r^{-1}=0$ and
$\Delta_\delta r^{-1-\gamma}
      =\gamma(1+\gamma)r^{-3-\gamma}$,
\[
 \Delta_g W\le-c_\gamma r^{-3-\gamma}
\]
for large $r_0$.  Choose a constant multiple of $W$ to dominate the
inner boundary values and the right side of
Equation~\ref{so:transformed-poisson}; this is possible because
$\gamma<\beta$.  The zero outer boundary values lie between the
positive and negative barriers.  The maximum principle therefore gives
$|\Phi(Z)|\le Cr^{-1}$, and the bounds for $\Phi'$ give
\begin{equation}\label{so:zero-end-constant}
 |Z|+|t|\le Cr^{-1}
\end{equation}
on every truncated end.  In particular no undetermined additive
constant is introduced in the exhaustion.

We record the derivative decay rather than infer it from an unscaled
estimate.  Rescale an annulus of radius $r$ to unit size.  The bounded
right side of Equation~\ref{so:transformed-poisson} and
Equation~\ref{so:zero-end-constant} give, by local Poisson $W^{2,s}$
estimates for $s>3$,
$|d\Phi(Z)|\le Cr^{-2}$ and a scaled H\"older bound for its gradient.
The exponentially small error has all unscaled derivatives
exponentially small, so it remains negligible after this rescaling.
The other right-side term is a smooth function of $Z$ times the
prescribed weight $\rho_0$; its scaled H\"older norms are bounded by
$Cr^{-3-\beta}$.  Schauder estimates give
$D^2Z=O(r^{-3})$.  Differentiating the transformed equation and
repeating the scaled estimates inductively gives
$Z=O_j(r^{-1})$ for every fixed $j$.  Since $t-Z$ is exponential,
the same holds for $t$.  Finally the untransformed equation gives
\[
 \Delta_g t=O(r^{-3-\beta})+O(|dZ|^2)+O(e^{-cr})
           =O(r^{-3-\beta}),
\]
because $\beta<1$ and $|dZ|^2=O(r^{-4})$.

Take any sequence $R\to\infty$.  The compact-patch estimates and a
diagonal subsequence give a smooth solution on $\Omega_N$, with all
bounds above and the stated end estimates.  Initially the limit has
$t\ge-\epsilon$.  If equality held at a point, it would be an interior
minimum of a smooth solution and the same strict contradiction in
Lemma~\ref{es:floor} would apply.  Hence $t>-\epsilon$.
No limit in $N$ has been taken.
Since $Z\to0$ and $N\epsilon>2$, the cutoff
$\vartheta(N(t+\epsilon))$ vanishes sufficiently far out.  The exact
zero-trace expression for $\mathcal T$ and the exponential trace error
give $\mathcal T=O(r^{-4})+O(e^{-cr})$ there.  Also
$\eta a_w|df|$ is exponential, $\rho=O(r^{-3-\beta})$, and
$u=e^Zl(Z)+O(e^{-cr})\to1$.  These functions are integrable over the
three-dimensional end, and smoothness handles the compact part.
Thus $u\Xi\in L^1(\Omega_N)$ and the final assertions follow.
\end{proof}
